\chapter{Khai báo trung thực học thuật}
\label{app:trung-thuc}

Phụ lục này khai báo phần tự làm, phần kế thừa, công cụ hỗ trợ, nguồn dữ liệu và những điều
báo cáo không khẳng định. Mục tiêu là để người đọc đánh giá đúng mức độ tin cậy của từng kết
quả.

\section{Phần tự cài đặt}
Nhóm tự viết:
\begin{itemize}
  \item Engine bitboard: biểu diễn 121 bit, \texttt{has\_five}, \texttt{candidates},
        \texttt{window\_counts}, \texttt{winning\_squares}, \texttt{evaluate}, Zobrist, đối
        xứng D4.
  \item Tám cấp tác tử: Random, Rule-based, Alpha-Beta, MCTS, Policy SL, REINFORCE, PUCT,
        vòng lặp AlphaZero.
  \item Arena: chạy trận, round-robin, khai cuộc, thống kê Wilson/bootstrap/Bradley-Terry,
        benchmark.
  \item Huấn luyện: sinh dữ liệu chuyên gia, SL, REINFORCE, vòng lặp tự đấu có gating.
  \item Hạ tầng: config tập trung, log JSONL, quản lý run, giao diện FastAPI, script sinh bảng
        và hình cho báo cáo.
\end{itemize}

\section{Phần kế thừa}
Thư viện và công cụ bên ngoài:
\begin{itemize}
  \item \textbf{NumPy, PyYAML, huggingface\_hub}: tính toán mảng, đọc cấu hình, tải model.
  \item \textbf{PyTorch}: mạng nơ-ron, huấn luyện, suy luận GPU.
  \item \textbf{FastAPI, Uvicorn}: giao diện web.
  \item \textbf{pandas, Matplotlib, Jupyter}: phân tích log, vẽ hình.
  \item \textbf{pytest, ruff}: kiểm thử và kiểm tra mã.
  \item \textbf{Cython}: lõi tăng tốc tùy chọn cho Alpha-Beta và PUCT.
\end{itemize}
Thuật toán tham khảo từ các công trình được trích dẫn ở Chương~\ref{ch:ly-thuyet}. Không có
đoạn mã nào được sao chép từ engine bên ngoài.

\section{Công cụ hỗ trợ}
Báo cáo và một phần mã được soạn với sự hỗ trợ của trợ lý lập trình AI (Claude Code). Phạm vi hỗ
trợ gồm: gợi ý cấu trúc báo cáo, soát lỗi chính tả, viết một số đoạn mô tả thuật toán, và hỗ
trợ kỹ thuật khi chạy thực nghiệm. Mọi số liệu trong báo cáo được sinh từ log chạy thật, không
do AI tạo ra. Nhóm chịu trách nhiệm kiểm tra lại mã, số liệu và kết luận. Nếu quy định môn học
yêu cầu khai báo chi tiết hơn, nhóm bổ sung ở bản nộp.

\section{Nguồn dữ liệu}
Mọi dữ liệu huấn luyện và đánh giá do nhóm tự sinh:
\begin{itemize}
  \item Dữ liệu chuyên gia: Alpha-Beta của nhóm, không dùng ván người hay engine ngoài.
  \item Checkpoint SL/RL: huấn luyện từ dữ liệu trên.
  \item Checkpoint L8: tự đấu từ L5, không dùng dữ liệu ngoài.
  \item Khai cuộc: sinh ngẫu nhiên có seed trong vùng 5×5 giữa bàn.
  \item Log đấu: do arena của nhóm ghi.
\end{itemize}
Không có dữ liệu cá nhân, dữ liệu nhạy cảm hay dữ liệu lấy từ internet.

\section{Điều báo cáo khẳng định và không khẳng định}
\begin{center}\small
\begin{tabularx}{\textwidth}{@{}p{.3\textwidth} X X@{}}
\toprule
\textbf{Vấn đề} & \textbf{Có bằng chứng} & \textbf{Không khẳng định} \\
\midrule
Sức mạnh tác tử & Điểm trong cấu hình, bộ khai cuộc, ngân sách đã ghi & Không khẳng định ``mạnh nhất nói chung'' \\
Tối ưu nước đi & Phân tích cơ hội thắng/chặn tức thì & Không khẳng định nước đi tối ưu dài hạn \\
Lợi thế đi trước & Đo trên từng cặp tác tử, từng bộ khai cuộc & Không khẳng định hằng số cho mọi cấu hình \\
AlphaZero-inspired & Vòng lặp, gating, checkpoint, kết quả & Không khẳng định tái hiện AlphaZero gốc \\
Cải thiện do RL & So sánh checkpoint, learning curve & Không tách hoàn toàn khỏi nhiễu seed/khai cuộc \\
Ablation & A1/A2/A3/A4/A5/A7 & A6/A7b/A8b/A5b chưa đủ, ghi rõ \\
\bottomrule
\end{tabularx}
\end{center}

\section{Lỗi đã phát hiện và sửa}
\paragraph{Quy ước dấu value.} Nhãn value trong dữ liệu SL v1/v2 là tuyệt đối (+1 nếu X thắng),
trong khi mạng và RL cần tương đối (người đi thắng). Lỗi làm baseline REINFORCE và value của
PUCT sai dấu ở khoảng một nửa số thế. Đã sửa ở SL v3 bằng \texttt{mover\_relative\_values}; đo
lại MSE 0,356 và tương quan +0,778. Các run SL v1/v2 vẫn giữ nhãn cũ và được ghi rõ trong báo
cáo.

\paragraph{Tham số model trong YAML.} File \texttt{resnet\_small.yaml} ghi \texttt{fc\_hidden:
64} nhưng code dùng 32. Báo cáo dùng code làm chuẩn và ghi chú chỗ lệch.

\paragraph{Đường dẫn input trong báo cáo.} Một số bảng ban đầu dùng đường dẫn tương đối sai
(\texttt{so-lieu-*.tex} thay vì \texttt{ban-mo-rong/so-lieu-*.tex}); đã sửa trước khi build
bản cuối.

\section{Những gì không nằm trong báo cáo}
Để tránh hiểu nhầm, các mục sau \emph{không} được coi là kết quả của báo cáo:
\begin{itemize}
  \item Chế độ cùng thời gian mỗi nước cho đủ tám tác tử (chưa chạy).
  \item Lặp nhiều seed (chưa chạy).
  \item Tập khai cuộc giữ riêng (chưa chạy).
  \item A6, A7b, A8b, A5b đầy đủ (chưa chạy hoặc chưa đủ mẫu).
  \item So sánh với engine bên ngoài như Rapfi (chưa chạy).
  \item Biến thể luật chặn hai đầu, Renju, bàn 15×15 (ngoài phạm vi).
\end{itemize}
Các mục này được liệt kê ở Phụ lục~\ref{app:cho-bo-sung} với công sức và tác động ước tính.

\section{Trách nhiệm}
Nhóm chịu trách nhiệm về tính chính xác của số liệu và mã nguồn. Nếu phát hiện sai sót, nhóm
sẽ sửa và ghi lại trong lịch sử commit; báo cáo không cố che giấu các run thất bại hoặc các
phần chưa làm. Cách trình bày ``những gì chưa làm'' là một phần của tiêu chuẩn trung thực học
thuật, không phải điểm yếu cần giấu.